\paragraph{Exact critical-cycle and finite-core linking input for D.2.45.}
This calculation concerns the primitive Mumford patterns
OS45/II with $(a,b)=(1,8)$ and zero index $2$ or $3$,
OS55/II with $(a,b)=(5,4)$ or $(1,20)$ and zero index $3$,
and OS55/default with $(a,b)=(1,20)$ and zero index
$2,3,4$. Indices are the zero-based indices of the retained
marked tuples. Each II cap may be original or a free good-II
logarithm. There are no additional regular logarithms.
The statement below supplies actual cycles and rational
linking, independently of any subsequent framed-bordism
or refined-Pontryagin interpretation.

Let $C$ denote the literal complex critical cycle of the
corrected differential. Every primitive rank-one Mumford
cap contributes its transverse elliptic curve and the
rank-two cusp contributes its three double curves,
with the unimodular Tate and support-function preimages
already established for D.2.44. An original II cap
contributes $2E_{II}$; a good-II logarithmic cap has no
critical contribution. Write
\[
 \Omega=a e_1^*e_2^*-b\delta^*\gamma^*,\qquad K=\Omega(C).
\]
Here each curve is first lifted as the actual integral
regular-fibre cycle, and only then evaluated by $\Omega$.
The exact values are
\[
\begin{array}{c|c|c|r|c}
\hbox{profile}&(a,b)&\hbox{zero and II choice}&K&
                   \operatorname{ord}[C]\text{ in }H_2\\\hline
45/II&(1,8)&2\text{ or }3,\ \mathrm{original}&-32&1\\
45/II&(1,8)&2\text{ or }3,\ \mathrm{log}&-16&1\\
55/II&(5,4)&3,\ \mathrm{original}&-108&1\\
55/II&(5,4)&3,\ \mathrm{log}&-100&1\\
55/II&(1,20)&3,\ \mathrm{original}&-68&5\\
55/II&(1,20)&3,\ \mathrm{log}&-28&5\\
55/\mathrm{default}&(1,20)&2&-108&5\\
55/\mathrm{default}&(1,20)&3\text{ or }4&-68&5.
\end{array}                                                   \tag{D.EC1}
\]
Thus the $c=20$ critical cycle is not integrally null.
All linking calculations below retain its exact order
$D=5$ and divide the intersection with a chain bounding
$DC$ by $D$.

\emph{The original-II curve and its marking.}
At a II fibre both sections are narrow. If $A$ is its
full marked matrix, put $B=A|_{\langle e_1,e_2\rangle}$,
$q=A_{\mathrm{ell},\gamma}$ and
\[
 z=(B-I)^{-1}q\in\mathbb Z^2,\qquad
                         \ell=(-z,0,1).
                                                               \tag{D.EC2}
\]
Here $\det(B-I)=1$, and $A\ell=\ell$ in the full
four-dimensional lattice. The original cap retracts
to the principal two-torus bundle over the cuspidal
rational curve. Its component Chern vector is zero.
The rigidified logarithm comparison used in D.2.31--32
identifies its principal circles with $\delta,\ell$,
including their boundary marking. The orbit of these
two circles through the cusp point is exactly the
reduced singular elliptic curve $E_{II}$. Their regular
torus specializes with degree one, so its literal
preimage is
\[
                             [E_{II}]=\delta\wedge\ell.
                                                               \tag{D.EC3}
\]
Adding a multiple of $\delta$ to the long lift has no
effect on (D.EC3). The complex torus orientation is
the positive short/long period orientation, the same
one used for the transverse curves of D.2.44.
In coordinates the reduced cusp is $q=y^2-x^3$.
Its critical ideal is $(x^2,y)$ and has transverse
length two. This proves the coefficient $2$ rather
than replacing the critical scheme by its reduced
support. For a good-II log, the unramified free
quotient of the good cover has regular corrected
differential, so this contribution is absent.

The actual vectors needed in (D.EC2) are particularly
simple:
\[
\begin{array}{c|c|c}
\hbox{profile and pair}&\ell&\psi_\delta\\\hline
45/II, (1,8)&(-8,0,0,1)&(0,-1,1,0)\\
55/II, (5,4)\text{ or }(1,20)&(0,0,0,1)&(0,0,1,0).
\end{array}                                                   \tag{D.EC4}
\]
The row $\psi_\delta$ is the integral projection onto
the principal $\delta$ circle, not an assumed coordinate
projection. To derive it in general, set
$t=A_{\delta,\mathrm{ell}}(B-I)^{-1}$. The invariant
elliptic graph is $(x,t x,0)$, and the decomposition
into this graph, $\delta$, and $\ell$ gives
\[
 \psi_\delta(\lambda)=\lambda_\delta-
                  t(\lambda_{\mathrm{ell}}+z\lambda_\gamma),
 \qquad \psi_\ell(\lambda)=\lambda_\gamma.
                                                               \tag{D.EC5}
\]
This proves the precise integral projection used in
the finite Thom pairings below.

\emph{A core three-manifold that is a full torsion generator.}
For a free II logarithm let
\[
 m=6k,\quad v=r\delta+n\ell,\quad
                   \sigma=n-me=\pm1 .                     \tag{D.EC6}
\]
These are all the simply connected one-II-log data of
D.2.30, with arbitrary $k>0,r,e$. The core surface
$D_i$ has central base lattice
\[
 \langle\delta,\ell,g\mid mg=r\delta+n\ell\rangle
                                                  \simeq\mathbb Z^2.
\]
Its elliptic monodromy is trivial on $\delta,\ell$
and is $B$ on $g$. The primitive character
\[
 \chi(\delta)=1,\qquad \chi(\ell)=-\sigma r,
                         \chi(g)=-\sigma re               \tag{D.EC7}
\]
defines a connected oriented three-manifold
$N=\chi^{-1}(\mathrm{point})\subset D_i$.
Its positive normal direction in the core may be
chosen along $\delta$, since $\chi(\delta)=1$.
The character respects the power relation because
$r-\sigma rn=-\sigma rem$.
Put $h=g+\sigma re\delta$ and
$k_0=\ell+\sigma r\delta$. They generate its kernel
with the single relation $mh=nk_0$.
The coprimality $\gcd(m,n)=1$ follows from
$n-me=\pm1$. Choose integers $u,v_0$ with
$un+v_0m=1$. Then $\kappa=uh+v_0k_0$ generates
$\ker\chi$, with $h=n\kappa$ and $k_0=m\kappa$.
Consequently $N$ is the elliptic mapping torus with
monodromy $B^u$. Since $6\mid m$ and
$n\equiv\sigma\pmod6$, one has $u\equiv\sigma\pmod6$;
the monodromy is exactly $B^\sigma$. In particular
$\det(B^\sigma-I)=1$ and $b_1(N;\mathbb F_2)=1$.
Reversing the mapping-torus endpoint convention
replaces $B^\sigma$ by its inverse and leaves these
statements unchanged.

The full normal Thom columns of this finite cap are
indexed by the four actual loops $\lambda\in\Lambda$
and the root loop $g$. Their intersections with $N$
are exactly
\[
 \chi(\lambda)=\psi_\delta(\lambda)
                              -\sigma r\psi_\ell(\lambda),
                 \qquad \chi(g)=-\sigma re.                \tag{D.EC8}
\]
Indeed a normal Thom three-chain is the normal disc
bundle over that core loop, and intersection with
the cooriented level set counts its character value.
There is no multiplication by the logarithmic order:
that multiplication already occurs in the norm part
of the boundary column. Formula (D.EC8) also retains
the actual elliptic-to-$\delta$ shear in (D.EC5).

\emph{A literal all-parameter chain relation.}
The integral fibre coinvariant for every row in (D.EC1)
is $\mathbb Z$, primitively measured by $\Omega$.
Let $s$ be a regular-fibre two-cycle with $\Omega(s)=1$.
All monodromies fix $\delta$. Each bar correction is
zero on the single loop $\delta$, including the
transported last boundary, and
\[
                         \sum_iJ_i\delta=0.               \tag{D.EC9}
\]
Take the finite normal Thom column for $\lambda=\delta$,
plus $m$ times the ordinary clutch $\delta$ column and
$m$ times every other filled-cap $\delta$ sweep.
The graph coordinates cancel by (D.EC9). Its literal
fibre part is
\[
                  -\delta\wedge v+me\,\delta\wedge\gamma,
 \qquad\Omega\bigl(-\delta\wedge v+me\delta\gamma\bigr)
                                      =b(n-me)=b\sigma .
                                                               \tag{D.EC10}
\]
All added nonfinite cap sweeps bound their actual
principal-circle orbit discs. If their integral fixed
basis does not explicitly contain $\delta$, changing
to one which does introduces only the retained fibre
monodromy relations; these changes are supported away
from the finite core $N$.

The difference between the fibre vector in (D.EC10)
and $b\sigma s$ is in the integral monodromy image,
not merely its saturation: its Smith factors are all
units and its primitive annihilator is $\Omega$.
These regular-piece chains also avoid $N$.
Thus (D.EC10), multiplied by $\sigma$, is an actual
chain bounding $bs$. Its intersection with $N$ is
\[
                              \sigma .                    \tag{D.EC11}
\]
This integer is a unit modulo $b$. Therefore $N$
is a full generator of $H_3(X)=\mathbb Z/b$, by
the nonsingular torsion linking pairing, including
when $b=20$. The fifth normal-column value in
(D.EC8) is essential for this stronger character;
it is not discarded when it is nonzero.

The difference $C-Ks$ also bounds a regular-piece
chain, after the specified local curve preimages
are used. Every local comparison chain is away
from $N$, since the II logarithmic cap has no
critical curve. If $D=\operatorname{ord}[C]$, the
actual chain for $DC$ therefore has intersection
\[
                       \sigma\frac{DK}{b}.
                                                               \tag{D.EC12}
\]
Equivalently its literal rational linking, with the
coorientations just specified, is
\[
                       \mathcal L(N,C)
                           =\frac{K}{b\sigma}.            \tag{D.EC13}
\]
This proves the formula for every integer parameter
in (D.EC6). In particular $DK/b$ equals $-2$ for
the II-log $c=8$ rows, $-25$ for the II-log $c=4$
row, and $-7$ for the II-log $c=20$ row. Their
two-primary parities are respectively even, odd,
odd, independently of every logarithmic order,
twist and ordinary clutch.

For an original II cap, or a semistable default
row, there is no finite log in this calculation.
The ordinary clutch has $e=\pm1$. The same sum
of all peripheral $\delta$ sweeps has fibre part
$e\delta\gamma$, with $\Omega$-value $-be$.
The ordinary regular three-torus
$T^3=\langle e_1,e_2,\gamma\rangle$ consequently
has linking $-e/b$ with $s$, so it is a full
$H_3$ generator. Its literal rational critical
linking is
\[
                           \mathcal L(T^3,C)=-eK/b.        \tag{D.EC14}
\]
For $c=20$ this calculation bounds $5C$ and then
divides by $5$; it does not replace $C$ by zero
in integral homology.

\emph{Exact retained certificates.}
The program \path{code/manifolds_extended_critical_cycles.py}
retains all curve preimages, full bounding coefficients,
the actual root lattice and characters, and both regular
and finite-core intersections. Its $158$ full Smith
replays comprise both original clutch signs and, for
each finite-II row, $k=1,2,3$, $r=0,1$, $e=-1,0,1$ and
both signs of $n-me$. They verify (D.EC12), as well as
the proper image of the regular $T^3$ in the logarithmic
cases; that regular torus is not silently called a
generator there.

More importantly, the $11$ constant certificates in
\path{records/manifolds_extended_critical_symbolic.json}
prove the unbounded identities. They store integral
monodromy preimages $c_1,c_2$ of
$-\delta\ell-bs$ and $\delta\gamma+bs$.
For arbitrary $m,n,e,r$ the literal identity is
\[
 R(nc_1+me c_2)=
 -\delta\wedge(r\delta+n\ell)+me\delta\gamma
                                  -b(n-me)s,
\]
where $R$ is the complete fibre monodromy matrix.
They also store the integral preimage of $C-Ks$,
the zero bar corrections and (D.EC9). Thus no
finite parameter search is used to infer the
all-order linking formula. The refined invariant
consequence requires the independently proved
ambient framing comparison; it is not inferred
from these chain calculations alone.
